% !TEX root = ../main.tex
\section{Methodology}
\label{sec:methodology}

The methodology followed in this work has five parts. \cref{sec:core-translation} establishes that the translated core is numerically faithful to PEANUTS and computationally competitive, the precondition for the rest of the chapter. \cref{sec:extension-methodology,sec:bsm-methodology} describe how that validated core is extended into a modular propagation framework and into beyond-the-Standard-Model scenarios. \cref{sec:detector-validation-methodology} then closes the loop by validating the full source--medium--detector--inference chain against real data, and \cref{sec:opensource-release} documents the release of the resulting code as a public, reusable package for the community.

\subsection{Core translation: methodology, validation and benchmarking}
\label{sec:core-translation}

\subsubsection{Translation methodology}

The development of \texttt{tpeanuts} follows a systematic translation process: each PEANUTS \parencite{Gonzalo2024} module (PMNS construction, vacuum and matter Hamiltonians, matter mixing, numerical propagation, solar production and Earth propagation, detector exposure) is first interpreted from the underlying theoretical formalism of \cref{sec:propagation-fundamentals}, separating the physical algorithm from implementation-specific detail, and then reformulated using native, batched PyTorch \parencite{Paszke2019} tensor operations, so that large grids of energies, baselines and production points are evaluated within a single vectorised call. PEANUTS's physical conventions, the PMNS phase, the neutrino/antineutrino sign, mass-eigenstate ordering, matter-potential normalisation, are explicitly verified at every step: intermediate quantities (mixing matrices, Hamiltonians, eigenvalues, matter potentials, oscillation phases) are compared directly against the original implementation before validating the final oscillation probabilities, since sign or unit-conversion errors can otherwise produce numerically stable but physically incorrect results.

Central to this translation is the batched PyTorch tensor formulation detailed architecturally in \cref{sec:pipelines}: every quantity, from PMNS matrices to full transition probabilities, carries an explicit batch dimension, and every operation on it is applied simultaneously across that dimension by ordinary broadcasting, so that a parameter scan, a detector-response folding over thousands of energy bins, or a gradient-based fit all reuse the same vectorised call rather than looping explicitly in Python, moving between CPU/GPU and single/double precision with no change to the physics code. The low-level physics primitives preserve the autograd graph by default; the fast-forward routines behind most results in \cref{sec:experimental-reproduction-results,sec:bsm-results} wrap this in \texttt{torch.no\_grad()} for speed, while the inference models of \cref{sec:detector-validation-methodology} build on the same differentiable primitives directly, so the same gradient constructing the oscillation Hamiltonian also drives the LBFGS fit and the Laplace/Fisher uncertainty estimate.

\subsubsection{Validation strategy}
\label{sec:validation-strategy}

Validation proceeds in two independent directions. \textbf{Intrinsically}, the perturbative/analytical and the purely numerical (matrix-exponential) evolutors of \cref{sec:propagation-fundamentals} solve the same propagation problem by independent strategies and are checked against each other, to a common tolerance, wherever both apply. \textbf{Extrinsically}, numerical, code-to-code agreement is verified hierarchically against two independent references:

\begin{itemize}
  \item low-level quantities, such as PMNS matrix elements, kinetic phases and matter potentials;
  \item intermediate physical outputs, such as effective mixing angles and eigenvalues in matter and mass-eigenstate weights for solar neutrinos;
  \item final oscillation probabilities, energy- and production-averaged and, where a detector model exists, folded into predicted event rates.
\end{itemize}

Since PEANUTS only implements solar-neutrino propagation, the comparison against it is restricted to that regime: vacuum propagation, the solar MSW effect, production-weighted survival probabilities, and Earth regeneration. \texttt{nuSQuIDS} \parencite{nuSQuIDS}, which solves the propagation problem independently for both the solar and the atmospheric regimes, cross-checks the full scope of the implementation instead and, relying on a different numerical method, also exposes discrepancies from conventions shared between PEANUTS and \texttt{tpeanuts} but absent from an independent solver. Agreement is verified to machine precision for vacuum and atmospheric propagation, where \texttt{tpeanuts} and the reference codes solve the same closed problem; for solar and solar-Earth-detection observables, the comparison against \texttt{nuSQuIDS} instead shows percent-to-ten-percent-level differences, attributable to genuine convention and approximation differences rather than floating-point or integration error. Probability conservation, positivity, and unitarity of the evolution operator are checked throughout.

\subsubsection{Benchmarking}
\label{sec:benchmarking-methodology}

Beyond physical correctness, an important objective is to evaluate computational performance, both \textbf{intrinsically} (CPU vs GPU, float32 vs float64, numerical versus perturvative, etc.) and \textbf{extrinsically} (against PEANUTS and \texttt{nuSQuIDS}), since efficient evaluation of large numbers of configurations is a principal motivation for the tensor-based design.

\textbf{Intrinsic benchmarking.} Executing the batched evaluation of \cref{sec:core-translation} on a GPU means thousands of arithmetic units working through the whole batch in parallel, rather than one core working sequentially through a list: launching a computation has a fixed cost independent of batch size, so small batches favour a sequential implementation while large batches amortise that cost and let GPU parallelism dominate, the key to reading the results of \cref{sec:benchmark-results}. Tensors are pre-allocated on the target device, GPU timings are taken only after explicit synchronisation, and warm-up iterations are excluded from the reported, repetition-averaged timings; both CPU and GPU are evaluated in single- and double-precision arithmetic, the latter for direct comparison with the reference codes.

\textbf{Extrinsic benchmarking.} The same oscillation parameters, density profiles, precision and scenarios are used across all three packages wherever the underlying calculation is directly comparable, vacuum and atmospheric propagation in particular. The suite varies the number of energy points, baselines/density samples, batch size, device, precision and oscillation scenario; the metric reported throughout this thesis is wall-clock speed-up relative to PEANUTS and \texttt{nuSQuIDS}. \texttt{nuSQuIDS} evaluates one solver object per energy/angle point, while PEANUTS already vectorises with NumPy over a single trajectory; \texttt{tpeanuts} batches across the full grid at once, so its relative advantage over \texttt{nuSQuIDS} grows steadily with grid size, while over PEANUTS it only emerges once the batch is large enough to amortise the fixed cost of dispatching PyTorch operations.

\subsection{Extension to a modular architecture: sources, media and detectors}
\label{sec:extension-methodology}

The validated core of \cref{sec:validation-strategy} is deliberately decoupled from any specific neutrino source, propagation medium or detector: the same Hamiltonian and evolution machinery serve solar, atmospheric and reactor sources; vacuum, solar and Earth matter media, with interchangeable density models; and the detector models of \cref{tab:detector-experiments} built on top of the resulting probabilities. The same engine could serve accelerator sources with no change to the propagation core. This separation is what allows the framework to extend beyond PEANUTS's original solar-only scope, to atmospheric-neutrino propagation through the Earth, alternative solar and Earth density models, and new detectors, without touching the propagation core already validated above.

\subsection{BSM extensions: NSI and sterile neutrinos}
\label{sec:bsm-methodology}

Beyond-the-Standard-Model scenarios are incorporated as modifications of the same Hamiltonian (\cref{sec:nsi-formalism,sec:sterile-formalism}), validated by a signature specific to each extension. For \textbf{NSI}, the benchmark is the LMA-Dark degeneracy \parencite{MirandaTortolaValle2004,Esteban2018}: the simultaneous transformation \(\theta_{12}\to\pi/2-\theta_{12}\), \(\varepsilon_{ee}\to-(2+\varepsilon_{ee})\) must reproduce a day--night-asymmetry sign flip consistent with the near-degenerate solution, exercising the sign dependence of the implementation. For \textbf{sterile neutrinos}, three presets (a null-mixing reference, a global-fit-motivated best fit, and a short-baseline-anomaly benchmark) are checked against analytic expectations: exact recovery of the Standard-Model limit at zero sterile mixing and unitarity of the enlarged \(4\times4\) mixing matrix to machine precision..

\subsection{Inference methodology: the full source--medium--detector chain}
\label{sec:detector-validation-methodology}

Validating the detector and inference layers exercises the full chain built in the sections above end to end, source flux, medium propagation, detector response and statistical inference, against real, published experimental data, since no independent reference implementation exists to compare against. Three complementary checks are applied, in full where the detector model is complete and in part otherwise, across the implemented experiments:

\begin{itemize}
  \item a \textbf{closure test}, a prediction-reproduction check against real published data, evaluating the full forward model at the experiment's own published best-fit parameters.
  \item an \textbf{absolute-normalisation audit}, when the closure test reveals a systematic offset, tracing it through every stage of the forward-model chain against independently published values to identify and quantify its dominant sources;
  \item a \textbf{parameter-recovery fit}, minimising the appropriate likelihood of \cref{sec:inference-formalism} with the quasi-Newton LBFGS optimiser \parencite{LiuNocedal1989}, comparing the fitted values against the experiment's own published best fit and against the global-fit value from NuFIT \parencite{EstebanNuFit2024}.
\end{itemize}

This methodology is applied in full to Daya Bay, the primary case study of \cref{sec:dayabay-case-study}; in summary form, with a parameter-recovery fit, to IceCube DeepCore and SNO phase~II; and only up to the implementation/unit-test stage, with no fit to real data yet reported, for Borexino and SNO phase~I (\cref{sec:other-experiments}).

\subsection{Open-source release: a public, documented software package}
\label{sec:opensource-release}

Releasing \texttt{tpeanuts} as public, open-source software was a deliberate part of this methodology of work. The full source code, together with the notebooks used for every validation, benchmark and analysis in this thesis, is publicly available on GitHub under the GNU General Public License v3 (GPLv3), the same licence as the original PEANUTS implementation it builds on, with an editable, development-mode install (\texttt{pip install -e .}) and a documented layout. Beyond the code itself, six companion \LaTeX{} documents provide module-by-module architectural and physics detail beyond the scope of this thesis, each with its own figures and, for the inference documents, its own real-data results. Full details, including the repository URL and the notebook directory tree, are given in \cref{sec:supplementary-material}.
